\documentclass[11pt,a4paper,oneside]{article}

% -------------------------------------------------
% Basic packages
% -------------------------------------------------
\usepackage[T1]{fontenc}
\usepackage[utf8]{inputenc}
\usepackage[spanish,es-noquoting,es-nodecimaldot]{babel}
\usepackage{lmodern}
\usepackage{microtype}
\usepackage[a4paper,margin=1.18in,headheight=15pt]{geometry}
\usepackage{amsmath,amssymb,amsfonts,amsthm,mathtools}
\usepackage{bm}
\usepackage{enumitem}
\usepackage{booktabs}
\usepackage{array}
\usepackage{xcolor}
\usepackage{hyperref}
\usepackage{setspace}
\usepackage{fancyhdr}
\usepackage{titlesec}
\usepackage[most]{tcolorbox}
\usepackage[capitalize,noabbrev]{cleveref}

\hypersetup{
  colorlinks=true,
  linkcolor=blue!50!black,
  citecolor=red!55!black,
  urlcolor=teal!60!black,
  pdftitle={Fundamentos de Econometría -- Hoja de ejercicios: MCO},
  pdfauthor={César Galindo}
}

% -------------------------------------------------
% Theorem-like environments
% -------------------------------------------------
\theoremstyle{plain}
\newtheorem{teorema}{Teorema}[section]
\newtheorem{proposicion}[teorema]{Proposición}
\newtheorem{lema}[teorema]{Lema}

\theoremstyle{definition}
\newtheorem{exercise}[teorema]{Ejercicio}

% -------------------------------------------------
% Fancy page and heading style
% -------------------------------------------------
\definecolor{econblue}{RGB}{25,62,114}
\definecolor{econgold}{RGB}{160,120,40}
\pagestyle{fancy}
\fancyhf{}
\fancyhead[L]{\textsc{Fundamentos de Econometría -- ECO515}}
\fancyhead[R]{\textsc{César Galindo}}
\fancyfoot[C]{\thepage}
\renewcommand{\headrulewidth}{0.35pt}
\renewcommand{\footrulewidth}{0pt}
\titleformat{\section}{\Large\bfseries\color{econblue}}{\thesection.}{0.65em}{}
\titleformat{\subsection}{\large\bfseries\color{econblue}}{\thesubsection.}{0.65em}{}
\titlespacing*{\section}{0pt}{1.9em}{1.05em}
\titlespacing*{\subsection}{0pt}{1.45em}{0.95em}
\tcbset{colback=white,colframe=econblue!65!black,boxrule=0.5pt,arc=1.5pt}
\setlength{\parskip}{0.85em}
\setstretch{1.13}
\setlength{\parindent}{0pt}

% -------------------------------------------------
% Pedagogical boxes
% -------------------------------------------------
\newtcolorbox{ideaBox}[1][]{enhanced,breakable,colback=blue!2,colframe=econblue,
  borderline west={1.5pt}{0pt}{econblue},title=Idea central,fonttitle=\bfseries,#1}
\newtcolorbox{summaryBox}[1][]{enhanced,breakable,colback=gray!4,colframe=black!55,
  borderline west={1.5pt}{0pt}{black!70},title=Recordatorio,fonttitle=\bfseries,#1}
\newtcolorbox{warningBox}[1][]{enhanced,breakable,colback=red!2,colframe=red!55!black,
  borderline west={1.5pt}{0pt}{red!75!black},title=Error frecuente,fonttitle=\bfseries,#1}
\newtcolorbox{hintBox}[1][]{enhanced,breakable,colback=green!2!white,colframe=green!40!black,
  borderline west={1.5pt}{0pt}{green!45!black},title=Sugerencia,fonttitle=\bfseries\color{green!40!black},#1}

\newcommand{\Cov}{\operatorname{Cov}}
\newcommand{\Var}{\operatorname{Var}}
\newcommand{\Corr}{\operatorname{Corr}}
\newcommand{\E}{\mathbb{E}}
\newcommand{\SCR}{\mathrm{SCR}}

\title{\vspace{-1.5cm}\textbf{\Large Fundamentos de Econometría}\\[2pt]
\large Hoja de ejercicios: Mínimos Cuadrados Ordinarios (MCO)}
\author{César Galindo\\ \small Universidad Panamericana \textperiodcentered\ Licenciatura en Finanzas \textperiodcentered\ ECO515}
\date{\small Otoño 2026}

\begin{document}
\maketitle
\vspace{-0.8cm}
\tableofcontents
\vspace{0.3cm}

\begin{ideaBox}
Conjunto de ejercicios de práctica sobre MCO en la regresión lineal
simple $y_i=\beta_0+\beta_1x_i+u_i$, complementario a las notas
\emph{Suma de variables aleatorias, varianza, covarianza, y la deducción
de los estimadores de MCO}. Se agrupan en cinco bloques: (I) álgebra y
deducción, (II) cálculo numérico, (III) interpretación y aplicaciones,
(IV) preguntas de opción múltiple estilo examen CFA, (V) un reto
opcional. Varios ejercicios incluyen solución completa; otros se dejan
propuestos, algunos con sugerencia.
\end{ideaBox}

\begin{summaryBox}
Para toda la hoja, $\bar x,\bar y$ son las medias muestrales,
$S_{xx}=\sum_i(x_i-\bar x)^2$, $S_{yy}=\sum_i(y_i-\bar y)^2$,
$S_{xy}=\sum_i(x_i-\bar x)(y_i-\bar y)$, y los estimadores de MCO son
\[
  \hat\beta_1 = \frac{S_{xy}}{S_{xx}}, \qquad
  \hat\beta_0 = \bar y - \hat\beta_1\bar x, \qquad
  \hat u_i = y_i-\hat\beta_0-\hat\beta_1x_i,
\]
con $\sum_i \hat u_i=0$ y $\sum_i x_i\hat u_i=0$ (ecuaciones normales).
\end{summaryBox}

\subsection*{Deducción vía derivadas parciales}

Un camino alternativo a las fórmulas de MCO, complementario al del
Ejercicio~\ref{ej:sin-calculo} (que no usa cálculo diferencial): aquí se
usa, y se muestra completo, paso a paso, con las \textbf{dos} formas
usuales de resolver el sistema $2\times2$ resultante --por la inversa de
la matriz y por la regla de Cramer-- para que se vea que dan exactamente
el mismo resultado.

\emph{Paso 1: la función a minimizar.}
\[
  \SCR(b_0,b_1)=\sum_{i=1}^n(y_i-b_0-b_1x_i)^2 .
\]

\emph{Paso 2: derivar cada término y sumar.} Para un término genérico,
por la regla de la cadena,
\[
  \frac{\partial}{\partial b_0}(y_i-b_0-b_1x_i)^2 = 2(y_i-b_0-b_1x_i)(-1),
  \qquad
  \frac{\partial}{\partial b_1}(y_i-b_0-b_1x_i)^2 = 2(y_i-b_0-b_1x_i)(-x_i).
\]
Sumando sobre $i=1,\dots,n$ (derivar bajo el signo de suma es válido
porque es una suma finita):
\[
  \frac{\partial \SCR}{\partial b_0}=-2\sum_i(y_i-b_0-b_1x_i), \qquad
  \frac{\partial \SCR}{\partial b_1}=-2\sum_ix_i(y_i-b_0-b_1x_i).
\]

\emph{Paso 3: condiciones de primer orden.} Igualando ambas derivadas a
cero en $(b_0,b_1)=(\hat\beta_0,\hat\beta_1)$ y dividiendo entre $-2$:
\[
  \sum_i(y_i-\hat\beta_0-\hat\beta_1x_i)=0, \qquad
  \sum_ix_i(y_i-\hat\beta_0-\hat\beta_1x_i)=0.
\]

\emph{Paso 4: distribuir la suma y agrupar las incógnitas a la
izquierda.} Distribuyendo cada suma y separando los términos con
$\hat\beta_0,\hat\beta_1$ de los que no los tienen:
\[
  \sum_iy_i-n\hat\beta_0-\hat\beta_1\sum_ix_i=0
  \;\Longrightarrow\;
  n\,\hat\beta_0 + \Big(\sum_ix_i\Big)\hat\beta_1 = \sum_iy_i,
\]
\[
  \sum_ix_iy_i-\hat\beta_0\sum_ix_i-\hat\beta_1\sum_ix_i^2=0
  \;\Longrightarrow\;
  \Big(\sum_ix_i\Big)\hat\beta_0+\Big(\sum_ix_i^2\Big)\hat\beta_1=\sum_ix_iy_i.
\]
Dos ecuaciones lineales en las dos incógnitas $\hat\beta_0,\hat\beta_1$.

\emph{Paso 5: forma matricial.} Con $\mathbf b=(\hat\beta_0,\hat\beta_1)^\top$,
el sistema del Paso 4 es $A\mathbf b=\mathbf c$ con
\[
  A=\begin{pmatrix} n & \displaystyle\sum_ix_i \\[4pt]
  \displaystyle\sum_ix_i & \displaystyle\sum_ix_i^2\end{pmatrix}, \qquad
  \mathbf c=\begin{pmatrix}\displaystyle\sum_iy_i\\[4pt] \displaystyle\sum_ix_iy_i\end{pmatrix}
\]
($A$ es simétrica: la entrada fuera de la diagonal, $\sum_ix_i$, es la
misma arriba a la derecha y abajo a la izquierda, porque viene de la
misma suma cruzada en ambas ecuaciones del Paso 4).

\emph{Paso 6: la identidad de suma de cuadrados, y $\det A$.} Antes de
calcular $\det A$ conviene tener a mano la identidad que conecta la suma
``cruda'' $\sum_ix_i^2$ con $S_{xx}$. Partiendo de la definición
$S_{xx}=\sum_i(x_i-\bar x)^2$ y expandiendo el cuadrado:
\[
  S_{xx}=\sum_i(x_i-\bar x)^2 = \sum_i\big(x_i^2-2x_i\bar x+\bar x^2\big)
  = \sum_ix_i^2 - 2\bar x\sum_ix_i + \sum_i\bar x^2 .
\]
El último término es $n\bar x^2$ ($\bar x^2$ sumado $n$ veces, una por
cada $i$), y como $\bar x=\frac1n\sum_ix_i$ se tiene
$\sum_ix_i=n\bar x$, así que el término de en medio es
$2\bar x(n\bar x)=2n\bar x^2$. Sustituyendo ambos:
\[
  S_{xx} = \sum_ix_i^2 - 2n\bar x^2+n\bar x^2 = \sum_ix_i^2-n\bar x^2 .
\]
Falta escribir $n\bar x^2$ en términos de $\sum_ix_i$: sustituyendo
$\bar x=\frac1n\sum_ix_i$,
\[
  n\bar x^2 = n\Big(\frac1n\sum_ix_i\Big)^{\!2}
  = n\cdot\frac{1}{n^2}\Big(\sum_ix_i\Big)^2
  = \frac{n}{n^2}\Big(\sum_ix_i\Big)^2
  = \frac1n\Big(\sum_ix_i\Big)^2
\]
(el factor $n/n^2$ se simplifica a $1/n$, cancelando una $n$ arriba y
abajo). Sustituyendo de vuelta:
\[
  S_{xx} = \sum_ix_i^2 - \frac{(\sum_ix_i)^2}{n}. \tag{$\star$}
\]
Ahora sí, $\det A$: multiplicando \emph{ambos lados} de ($\star$) por
$n$,
\[
  n\,S_{xx} = n\sum_ix_i^2 - n\cdot\frac{(\sum_ix_i)^2}{n}
  = n\sum_ix_i^2-\frac{n}{n}\Big(\sum_ix_i\Big)^2
  = n\sum_ix_i^2-\Big(\sum_ix_i\Big)^2
\]
(de nuevo $n/n=1$, así que ese factor simplemente desaparece). El lado
derecho es exactamente $\det A = ps-qr$ para las entradas de $A$ del
Paso 5 ($p=n$, $s=\sum_ix_i^2$, $q=r=\sum_ix_i$), así que
\[
  \det A = n\sum_ix_i^2-\Big(\sum_ix_i\Big)^2 = n\,S_{xx},
\]
positivo siempre que $S_{xx}>0$ (es decir, que $x_i$ no sea constante
--la misma condición que hace a $A$ invertible). A partir de aquí, dos
caminos igual de válidos para resolver $A\mathbf b=\mathbf c$.

\paragraph{Camino 1: por la inversa de la matriz.}

\emph{Paso 7a: invertir $A$.} Con la fórmula de la inversa de una matriz
$2\times2$,
$\begin{psmallmatrix}p&q\\r&s\end{psmallmatrix}^{-1}
=\dfrac1{ps-qr}\begin{psmallmatrix}s&-q\\-r&p\end{psmallmatrix}$, y
usando $\det A=nS_{xx}$ del Paso 6:
\[
  A^{-1}=\frac1{n\,S_{xx}}
  \begin{pmatrix}\displaystyle\sum_ix_i^2 & -\displaystyle\sum_ix_i \\[4pt]
  -\displaystyle\sum_ix_i & n\end{pmatrix}.
\]

\emph{Paso 8a: $\mathbf b=A^{-1}\mathbf c$.} Multiplicando la matriz por
el vector $\mathbf c$ del Paso 5, entrada por entrada:
\[
  \begin{pmatrix}\hat\beta_0\\\hat\beta_1\end{pmatrix}=A^{-1}\mathbf c
  =\frac1{n\,S_{xx}}
  \begin{pmatrix}\Big(\sum_ix_i^2\Big)\Big(\sum_iy_i\Big)-\Big(\sum_ix_i\Big)\Big(\sum_ix_iy_i\Big)\\[6pt]
  n\displaystyle\sum_ix_iy_i-\Big(\sum_ix_i\Big)\Big(\sum_iy_i\Big)\end{pmatrix}.
\]

\emph{Paso 9a: simplificar la segunda entrada, $\hat\beta_1$.} Se
necesita la identidad análoga a ($\star$) del Paso 6, ahora para
$S_{xy}$. Partiendo de $S_{xy}=\sum_i(x_i-\bar x)(y_i-\bar y)$ y
expandiendo el producto:
\[
  S_{xy} = \sum_i\big(x_iy_i-x_i\bar y-\bar xy_i+\bar x\bar y\big)
  = \sum_ix_iy_i - \bar y\sum_ix_i - \bar x\sum_iy_i + n\bar x\bar y .
\]
Con $\sum_ix_i=n\bar x$ y $\sum_iy_i=n\bar y$, los dos términos de en
medio son $\bar y(n\bar x)=n\bar x\bar y$ y $\bar x(n\bar y)=n\bar x\bar y$;
sumados al último término, $-n\bar x\bar y-n\bar x\bar y+n\bar x\bar y=-n\bar x\bar y$, así que
\[
  S_{xy} = \sum_ix_iy_i - n\bar x\bar y .
\]
Sustituyendo $\bar x=\frac1n\sum_ix_i$, $\bar y=\frac1n\sum_iy_i$:
\[
  n\bar x\bar y = n\Big(\frac1n\sum_ix_i\Big)\Big(\frac1n\sum_iy_i\Big)
  = \frac{n}{n^2}\Big(\sum_ix_i\Big)\Big(\sum_iy_i\Big)
  = \frac1n\Big(\sum_ix_i\Big)\Big(\sum_iy_i\Big)
\]
(otra vez $n/n^2=1/n$), de modo que
\[
  S_{xy} = \sum_ix_iy_i-\frac{(\sum_ix_i)(\sum_iy_i)}{n}. \tag{$\star\star$}
\]
Multiplicando \emph{ambos lados} de ($\star\star$) por $n$:
\[
  n\,S_{xy} = n\sum_ix_iy_i - n\cdot\frac{(\sum_ix_i)(\sum_iy_i)}{n}
  = n\sum_ix_iy_i-\frac{n}{n}\Big(\sum_ix_i\Big)\Big(\sum_iy_i\Big)
  = n\sum_ix_iy_i-\Big(\sum_ix_i\Big)\Big(\sum_iy_i\Big)
\]
(de nuevo $n/n=1$). El lado derecho es exactamente el numerador de
$\hat\beta_1$ del Paso 8a, así que ese numerador es $n\,S_{xy}$, y
\[
  \hat\beta_1=\frac{n\,S_{xy}}{n\,S_{xx}}
  = \frac{n}{n}\cdot\frac{S_{xy}}{S_{xx}}
  = \frac{S_{xy}}{S_{xx}}
\]
(el factor $n/n=1$ se cancela una última vez, dejando el cociente
familiar).

\emph{Paso 10a: simplificar la primera entrada, $\hat\beta_0$.} Se suma
y se resta $(\sum_ix_i)^2(\sum_iy_i)/n$ dentro del numerador, y se usan
las mismas definiciones de $S_{xx}$ y $S_{xy}$ de los Pasos 6 y 9a:
\begin{align*}
  \Big(\sum_iy_i\Big)S_{xx} - S_{xy}\Big(\sum_ix_i\Big)
  &= \Big(\sum_iy_i\Big)\Big[\sum_ix_i^2-\frac{(\sum_ix_i)^2}{n}\Big]
   - \Big[\sum_ix_iy_i-\frac{(\sum_ix_i)(\sum_iy_i)}{n}\Big]\Big(\sum_ix_i\Big) \\
  &= \Big(\sum_iy_i\Big)\Big(\sum_ix_i^2\Big) - \Big(\sum_ix_i\Big)\Big(\sum_ix_iy_i\Big)
   \underbrace{-\,\frac{(\sum_ix_i)^2(\sum_iy_i)}{n}+\frac{(\sum_ix_i)^2(\sum_iy_i)}{n}}_{=\,0} \\
  &= \Big(\sum_ix_i^2\Big)\Big(\sum_iy_i\Big)-\Big(\sum_ix_i\Big)\Big(\sum_ix_iy_i\Big),
\end{align*}
es decir, el numerador de la primera entrada del Paso 8a es exactamente
$(\sum_iy_i)S_{xx}-S_{xy}(\sum_ix_i)$. Dividiendo entre $\det A=n\,S_{xx}$
y separando la fracción en dos antes de simplificar:
\[
  \hat\beta_0=\frac{\big(\sum_iy_i\big)S_{xx}-S_{xy}\big(\sum_ix_i\big)}{n\,S_{xx}}
  =\frac{\big(\sum_iy_i\big)S_{xx}}{n\,S_{xx}}-\frac{S_{xy}\big(\sum_ix_i\big)}{n\,S_{xx}} .
\]
En el primer término $S_{xx}/S_{xx}=1$ se cancela, dejando
$(\sum_iy_i)/n=\bar y$; en el segundo se reagrupa
$S_{xy}/S_{xx}=\hat\beta_1$ y queda $(\sum_ix_i)/n=\bar x$:
\[
  \hat\beta_0 = \bar y-\hat\beta_1\bar x .
\]

\paragraph{Camino 2: por la regla de Cramer.}

\emph{Paso 7b: enunciado de Cramer.} Para $A\mathbf b=\mathbf c$ con
$\det A\neq0$,
\[
  \hat\beta_0=\frac{\det A_0}{\det A}, \qquad \hat\beta_1=\frac{\det A_1}{\det A},
\]
donde $A_0$ es $A$ con su \emph{primera} columna (los coeficientes de
$\hat\beta_0$) reemplazada por $\mathbf c$, y $A_1$ es $A$ con su
\emph{segunda} columna reemplazada por $\mathbf c$:
\[
  A_0=\begin{pmatrix}\displaystyle\sum_iy_i & \displaystyle\sum_ix_i \\[4pt]
  \displaystyle\sum_ix_iy_i & \displaystyle\sum_ix_i^2\end{pmatrix}, \qquad
  A_1=\begin{pmatrix} n & \displaystyle\sum_iy_i \\[4pt]
  \displaystyle\sum_ix_i & \displaystyle\sum_ix_iy_i\end{pmatrix}.
\]

\emph{Paso 8b: $\det A_0$ y $\det A_1$.} Desarrollando cada determinante
$2\times2$ por la fórmula usual ($ps-qr$):
\[
  \det A_0 = \Big(\sum_iy_i\Big)\Big(\sum_ix_i^2\Big) - \Big(\sum_ix_i\Big)\Big(\sum_ix_iy_i\Big),
  \qquad
  \det A_1 = n\sum_ix_iy_i - \Big(\sum_ix_i\Big)\Big(\sum_iy_i\Big).
\]
Éstas son, entrada por entrada, las mismas dos expresiones que aparecen
en el numerador de $\mathbf b=A^{-1}\mathbf c$ del Paso 8a --no es
casualidad: para una matriz $2\times2$, la fórmula de la inversa y la
regla de Cramer son álgebraicamente la misma construcción, sólo
organizada de otra manera.

\emph{Paso 9b: dividir entre $\det A$.} Del Paso 9a, $\det A_1=n\,S_{xy}$;
del Paso 10a, $\det A_0=(\sum_iy_i)S_{xx}-S_{xy}(\sum_ix_i)$. Dividiendo
cada uno entre $\det A=n\,S_{xx}$ (Paso 6):
\[
  \hat\beta_1=\frac{\det A_1}{\det A}=\frac{n\,S_{xy}}{n\,S_{xx}}
  =\frac{n}{n}\cdot\frac{S_{xy}}{S_{xx}}=\frac{S_{xy}}{S_{xx}}
\]
(el mismo factor $n/n=1$ del Paso 9a). Para $\hat\beta_0$, se separa la
fracción en dos antes de simplificar:
\[
  \hat\beta_0=\frac{\det A_0}{\det A}
  =\frac{\big(\sum_iy_i\big)S_{xx}-S_{xy}\big(\sum_ix_i\big)}{n\,S_{xx}}
  =\frac{\big(\sum_iy_i\big)S_{xx}}{n\,S_{xx}}-\frac{S_{xy}\big(\sum_ix_i\big)}{n\,S_{xx}} .
\]
En el primer término, $S_{xx}/S_{xx}=1$ se cancela, dejando
$(\sum_iy_i)/n=\bar y$; en el segundo, se reagrupa
$S_{xy}/S_{xx}=\hat\beta_1$ y queda $(\sum_ix_i)/n=\bar x$:
\[
  \hat\beta_0 = \bar y - \hat\beta_1\,\bar x .
\]

\emph{Conclusión.} Los dos caminos --inversa de la matriz y regla de
Cramer-- recuperan exactamente $\hat\beta_1=S_{xy}/S_{xx}$ y
$\hat\beta_0=\bar y-\hat\beta_1\bar x$, las mismas fórmulas que el
Ejercicio~\ref{ej:sin-calculo} obtiene completando el cuadrado, sin usar
cálculo. Los tres caminos son válidos exactamente bajo la misma
condición, $S_{xx}>0$ (equivalente a $\det A\neq0$): la elección entre
ellos es de conveniencia algebraica, no de contenido matemático
distinto.

% ================================================================
\section{Álgebra y deducción}
% ================================================================

\begin{exercise}[Deducción sin cálculo diferencial]\label{ej:sin-calculo}
Pruebe, \textbf{sin usar derivadas}, que $(\hat\beta_0,\hat\beta_1)$
minimiza $\SCR(b_0,b_1)=\sum_i(y_i-b_0-b_1x_i)^2$, completando cuadrados
en dos pasos:
\begin{enumerate}[nosep,leftmargin=1.8em,label=(\alph*)]
  \item Para $b_1$ fijo, muestre que $b_0=\overline{(y-b_1x)}$ minimiza
  $\sum_i(y_i-b_1x_i-b_0)^2$, usando la identidad
  $\sum_i(a_i-b_0)^2 = \sum_i(a_i-\bar a)^2+n(\bar a-b_0)^2$ con
  $a_i=y_i-b_1x_i$.
  \item Sustituya ese $b_0$ óptimo de vuelta en $\SCR$ para obtener una
  función \emph{sólo} de $b_1$, exprésela como
  $S_{xx}\,b_1^2-2S_{xy}\,b_1+S_{yy}$, y complete el cuadrado en $b_1$
  para mostrar que se minimiza en $b_1=S_{xy}/S_{xx}$.
\end{enumerate}
\end{exercise}
\begin{proof}[Solución]
(a) $\sum_i(a_i-b_0)^2 = \sum_i\big[(a_i-\bar a)+(\bar a-b_0)\big]^2
= \sum_i(a_i-\bar a)^2 + 2(\bar a-b_0)\sum_i(a_i-\bar a) + n(\bar a-b_0)^2$.
El término cruzado se anula porque $\sum_i(a_i-\bar a)=0$ (suma de
desviaciones respecto a la media). Queda
$\sum_i(a_i-\bar a)^2+n(\bar a-b_0)^2$, suma de un término que no
depende de $b_0$ más un cuadrado no negativo que se anula exactamente en
$b_0=\bar a=\overline{(y-b_1x)}=\bar y-b_1\bar x$.

(b) Sustituyendo $b_0=\bar y-b_1\bar x$, el residual se vuelve
$y_i-b_0-b_1x_i = (y_i-\bar y)-b_1(x_i-\bar x)$, así que
\begin{align*}
  \SCR(b_1) &= \sum_i\big[(y_i-\bar y)-b_1(x_i-\bar x)\big]^2 \\
  &= \sum_i(y_i-\bar y)^2 - 2b_1\sum_i(x_i-\bar x)(y_i-\bar y)
     + b_1^2\sum_i(x_i-\bar x)^2 \\
  &= S_{yy} - 2S_{xy}\,b_1 + S_{xx}\,b_1^2 .
\end{align*}
Completando el cuadrado en $b_1$ (posible porque $S_{xx}>0$):
\[
  \SCR(b_1) = S_{xx}\left(b_1-\frac{S_{xy}}{S_{xx}}\right)^{\!2}
  + \left(S_{yy}-\frac{S_{xy}^2}{S_{xx}}\right).
\]
El segundo término no depende de $b_1$; el primero es
$S_{xx}$ (positivo) por un cuadrado, así que es mínimo exactamente
cuando el cuadrado se anula: $b_1=S_{xy}/S_{xx}=\hat\beta_1$. Esto
recupera las fórmulas de los estimadores de MCO (Recordatorio anterior)
sin usar cálculo diferencial en ningún paso --sólo la misma identidad de
``completar el cuadrado'' usada para probar la desigualdad de
Cauchy--Schwarz.
\end{proof}

\begin{exercise}[Invarianza ante un cambio de escala en $x$]\label{ej:invarianza-x}
Sea $x_i' = a+bx_i$ ($b\neq0$) un reescalamiento del regresor (por
ejemplo, pasar de tasas en decimales a tasas en puntos porcentuales).
Sean $\hat\beta_0,\hat\beta_1$ los estimadores de MCO de $y$ sobre $x$, y
$\hat\beta_0',\hat\beta_1'$ los de $y$ sobre $x'$. Demuestre que
\[
  \hat\beta_1' = \frac{\hat\beta_1}{b}, \qquad
  \hat\beta_0' = \hat\beta_0 - \hat\beta_1\,\frac{a}{b}.
\]
\end{exercise}
\begin{proof}[Solución]
Como $\bar x'=a+b\bar x$, se tiene $x_i'-\bar x' = b(x_i-\bar x)$. Por
tanto
\[
  S_{x'y} = \sum_i (x_i'-\bar x')(y_i-\bar y) = b\sum_i(x_i-\bar x)(y_i-\bar y) = b\,S_{xy},
\]
\[
  S_{x'x'} = \sum_i (x_i'-\bar x')^2 = b^2\sum_i(x_i-\bar x)^2 = b^2 S_{xx}.
\]
Entonces $\hat\beta_1' = S_{x'y}/S_{x'x'} = bS_{xy}/(b^2S_{xx}) =
\hat\beta_1/b$. Para el intercepto,
$\hat\beta_0'=\bar y-\hat\beta_1'\bar x' = \bar y -
\dfrac{\hat\beta_1}{b}(a+b\bar x) = \bar y - \hat\beta_1\bar x -
\hat\beta_1\dfrac{a}{b} = \hat\beta_0 - \hat\beta_1\dfrac{a}{b}$.
\end{proof}

\begin{exercise}[Invarianza ante un cambio de escala en $y$]\label{ej:invarianza-y}
Sea $y_i'=cy_i$ ($c\neq0$; por ejemplo, pasar de pesos a miles de pesos).
Demuestre que los estimadores de MCO de $y'$ sobre el mismo $x$ satisfacen
$\hat\beta_1'=c\,\hat\beta_1$ y $\hat\beta_0'=c\,\hat\beta_0$. (Imite el
argumento del Ejercicio~\ref{ej:invarianza-x}.)
\end{exercise}
\begin{proof}[Solución]
Como $\bar y'=c\bar y$, se tiene $y_i'-\bar y'=c(y_i-\bar y)$, mientras
que $x_i-\bar x$ no cambia. Por tanto
\[
  S_{xy'} = \sum_i(x_i-\bar x)(y_i'-\bar y') = c\sum_i(x_i-\bar x)(y_i-\bar y) = c\,S_{xy},
\]
y $S_{xx}$ no cambia. Entonces $\hat\beta_1' = S_{xy'}/S_{xx} = c\,S_{xy}/S_{xx} = c\,\hat\beta_1$,
y $\hat\beta_0' = \bar y'-\hat\beta_1'\bar x = c\bar y - c\hat\beta_1\bar x
= c(\bar y-\hat\beta_1\bar x) = c\,\hat\beta_0$.
\end{proof}

\begin{exercise}[Cambio de escala simultáneo en $x$ y $y$]\label{ej:invarianza-combinada}
Combine los Ejercicios~\ref{ej:invarianza-x} y \ref{ej:invarianza-y} para
mostrar que, si $x_i'=a+bx_i$ y $y_i'=cy_i$ simultáneamente
($b,c\neq0$), entonces
\[
  \hat\beta_1' = \frac{c}{b}\,\hat\beta_1, \qquad
  \hat\beta_0' = c\,\hat\beta_0 - \frac{ac}{b}\,\hat\beta_1
  = c\,\hat\beta_0 - a\,\hat\beta_1' .
\]
¿Qué le dice esto sobre por qué, al reportar una regresión, siempre hay
que especificar las unidades de $x$ y $y$?
\end{exercise}
\begin{proof}[Solución]
Como $x_i'-\bar x'=b(x_i-\bar x)$ y $y_i'-\bar y'=c(y_i-\bar y)$,
\[
  S_{x'y'} = \sum_i(x_i'-\bar x')(y_i'-\bar y') = bc\,S_{xy}, \qquad
  S_{x'x'} = \sum_i(x_i'-\bar x')^2 = b^2 S_{xx},
\]
así que
$\hat\beta_1' = S_{x'y'}/S_{x'x'} = bc\,S_{xy}/(b^2S_{xx}) =
(c/b)\,S_{xy}/S_{xx} = (c/b)\,\hat\beta_1$. Para el intercepto, con
$\bar x'=a+b\bar x$ y $\bar y'=c\bar y$,
\[
  \hat\beta_0' = \bar y'-\hat\beta_1'\bar x'
  = c\bar y - \frac{c}{b}\hat\beta_1(a+b\bar x)
  = c\bar y - \frac{ac}{b}\hat\beta_1 - c\hat\beta_1\bar x
  = c(\bar y-\hat\beta_1\bar x) - \frac{ac}{b}\hat\beta_1
  = c\,\hat\beta_0 - \frac{ac}{b}\hat\beta_1,
\]
que, usando $\hat\beta_1'=(c/b)\hat\beta_1$, también se escribe
$\hat\beta_0'=c\hat\beta_0-a\hat\beta_1'$. En palabras: reescalar $x$ y
$y$ no cambia la forma de la recta ajustada, sólo sus coeficientes --por
eso un $\hat\beta_1$ reportado sin unidades es ambiguo: el mismo
fenómeno económico puede dar una pendiente de $0.7$ o de $700$ según se
mida $x$ en decimales o en puntos base.
\end{proof}

\begin{exercise}[Residuales y valores ajustados no están correlacionados]\label{ej:ortog-fitted}
Demuestre que $\sum_i \hat u_i\,\hat y_i = 0$, donde
$\hat y_i=\hat\beta_0+\hat\beta_1x_i$ son los valores ajustados.
\end{exercise}
\begin{proof}[Solución]
\[
  \sum_i \hat u_i\hat y_i = \sum_i \hat u_i(\hat\beta_0+\hat\beta_1x_i)
  = \hat\beta_0\sum_i\hat u_i + \hat\beta_1\sum_i x_i\hat u_i
  = \hat\beta_0(0)+\hat\beta_1(0) = 0,
\]
usando las dos ecuaciones normales. Junto con $\sum_i\hat u_i=0$, esto
dice que el vector de residuales es \emph{ortogonal}, en el sentido
usual del producto punto, tanto al vector constante como al vector de
valores ajustados --la razón geométrica detrás de la descomposición
$\mathrm{SST}=\mathrm{SSE}+\mathrm{SSR}$ del
Ejercicio~\ref{ej:r2-rho2}.
\end{proof}

% ================================================================
\section{Cálculo numérico}
% ================================================================

\begin{exercise}\label{ej:numerico-grande}
Con la muestra $n=6$:
\begin{center}
\begin{tabular}{@{}ccccccc@{}}
\toprule
$i$ & 1 & 2 & 3 & 4 & 5 & 6 \\
\midrule
$x_i$ & 2 & 4 & 6 & 8 & 10 & 12 \\
$y_i$ & 3 & 5 & 6 & 8 & 7 & 10 \\
\bottomrule
\end{tabular}
\end{center}
calcule $\bar x,\bar y,S_{xy},S_{xx},S_{yy}$, los estimadores
$\hat\beta_0,\hat\beta_1$, y la suma de cuadrados residuales en el
mínimo, $\SCR_{\min}=S_{yy}-S_{xy}^2/S_{xx}$.
\end{exercise}
\begin{proof}[Solución]
$\bar x=42/6=7$, $\bar y=39/6=6.5$.
\begin{center}
\begin{tabular}{@{}crrrr@{}}
\toprule
$i$ & $x_i-\bar x$ & $y_i-\bar y$ & producto & $(x_i-\bar x)^2$ \\
\midrule
1 & $-5$ & $-3.5$ & $17.5$ & $25$ \\
2 & $-3$ & $-1.5$ & $4.5$  & $9$  \\
3 & $-1$ & $-0.5$ & $0.5$  & $1$  \\
4 & $1$  & $1.5$  & $1.5$  & $1$  \\
5 & $3$  & $0.5$  & $1.5$  & $9$  \\
6 & $5$  & $3.5$  & $17.5$ & $25$ \\
\midrule
$\sum$ & & & $S_{xy}=43$ & $S_{xx}=70$ \\
\bottomrule
\end{tabular}
\end{center}
$S_{yy}=3.5^2+1.5^2+0.5^2+1.5^2+0.5^2+3.5^2 = 12.25+2.25+0.25+2.25+0.25+12.25=29.5$.
\[
  \hat\beta_1 = \frac{43}{70} \approx 0.6143, \qquad
  \hat\beta_0 = 6.5 - 0.6143(7) = 6.5-4.3 = 2.2,
\]
donde $0.6143\times7=43\times7/70=301/70=4.3$ exactamente. Finalmente,
\[
  \SCR_{\min} = S_{yy}-\frac{S_{xy}^2}{S_{xx}} = 29.5-\frac{43^2}{70}
  = 29.5-\frac{1849}{70} \approx 29.5-26.414 \approx 3.086.
\]
\end{proof}

\begin{exercise}[Sensibilidad a un dato atípico]\label{ej:outlier}
Recuerde el ejemplo de las notas con $n=5$: $x=(1,2,3,4,5)$,
$y=(2,3,5,4,6)$, para el cual $\hat\beta_1=0.9$, $\hat\beta_0=1.3$.
Sustituya el último punto por $(5,20)$ --un dato atípico-- y recalcule
$\hat\beta_0,\hat\beta_1$. Compare con los valores originales y comente,
en una o dos líneas, qué tan sensible es MCO a un solo dato extremo.
\end{exercise}
\begin{proof}[Solución]
Con $y=(2,3,5,4,20)$: $\bar y=34/5=6.8$ ($\bar x=3$ no cambia, así que
$S_{xx}=10$ sigue igual). Las desviaciones de $y$ son
$(-4.8,-3.8,-1.8,-2.8,13.2)$, y con $x-\bar x=(-2,-1,0,1,2)$ los
productos son $9.6,\,3.8,\,0,\,-2.8,\,26.4$, con suma
$S_{xy}=37$. Entonces
\[
  \hat\beta_1 = \frac{37}{10}=3.7, \qquad
  \hat\beta_0 = 6.8-3.7(3) = 6.8-11.1=-4.3 .
\]
La pendiente pasó de $0.9$ a $3.7$ (más de cuatro veces mayor) y el
intercepto de $1.3$ a $-4.3$: un único punto atípico, con la misma $x$
que los demás, basta para dominar por completo la recta ajustada. Esto
ilustra por qué MCO no es robusto a valores atípicos --la función
objetivo $\SCR$ penaliza el error al \emph{cuadrado}, así que una
observación lejana pesa desproporcionadamente en la suma.
\end{proof}

\begin{exercise}[Practicar con desviaciones ya calculadas]\label{ej:numerico-desviaciones}
Para una muestra de $n=8$ se reportan únicamente las desviaciones
respecto a la media: $\sum_i(x_i-\bar x)(y_i-\bar y)=54$,
$\sum_i(x_i-\bar x)^2=36$, $\sum_i(y_i-\bar y)^2=100$, $\bar x=12$,
$\bar y=45$. Sin conocer los datos individuales, calcule
$\hat\beta_0,\hat\beta_1$, la correlación muestral $\hat\rho_{xy}$, y
$R^2$.
\end{exercise}
\begin{proof}[Solución]
$\hat\beta_1=54/36=1.5$. $\hat\beta_0=45-1.5(12)=45-18=27$.
$\hat\rho_{xy}=54/\sqrt{36\times100}=54/60=0.9$. Anticipando el
Ejercicio~\ref{ej:r2-rho2}, $R^2=\hat\rho_{xy}^2=0.81$: el regresor
explica el 81\% de la variación muestral de $y$. Nótese que las tres
cantidades se obtienen sólo de $S_{xy},S_{xx},S_{yy},\bar x,\bar y$, sin
necesidad de los datos individuales --exactamente el mismo principio
detrás de reportar una regresión mediante su tabla de resultados en vez
de la base de datos completa.
\end{proof}

% ================================================================
\section{Interpretación y aplicaciones}
% ================================================================

\begin{exercise}[$R^2$ es la correlación muestral al cuadrado]\label{ej:r2-rho2}
Defina $\mathrm{SST}=S_{yy}$, $\mathrm{SSE}=\sum_i(\hat y_i-\bar y)^2$
(suma de cuadrados explicada) y $\mathrm{SSR}=\sum_i\hat u_i^2$ (suma de
cuadrados residual), con $\hat y_i=\hat\beta_0+\hat\beta_1x_i$.
\begin{enumerate}[nosep,leftmargin=1.8em,label=(\alph*)]
  \item Pruebe la descomposición $\mathrm{SST}=\mathrm{SSE}+\mathrm{SSR}$,
  mostrando que el término cruzado en
  $\sum_i\big[(\hat y_i-\bar y)+\hat u_i\big]^2$ se anula usando las
  ecuaciones normales.
  \item Pruebe que $\mathrm{SSE}=S_{xy}^2/S_{xx}$.
  \item Concluya que $R^2:=\mathrm{SSE}/\mathrm{SST} = \hat\rho_{xy}^2$,
  donde $\hat\rho_{xy}=S_{xy}/\sqrt{S_{xx}S_{yy}}$ es la correlación
  muestral.
\end{enumerate}
\end{exercise}
\begin{proof}[Solución]
(a) $y_i-\bar y = (\hat y_i-\bar y)+\hat u_i$, así que
\[
  \mathrm{SST}=\sum_i(y_i-\bar y)^2
  = \sum_i(\hat y_i-\bar y)^2 + 2\sum_i(\hat y_i-\bar y)\hat u_i + \sum_i\hat u_i^2.
\]
El término cruzado es
$\sum_i(\hat y_i-\bar y)\hat u_i = \hat\beta_1\sum_i(x_i-\bar x)\hat u_i$
(pues $\hat y_i-\bar y=\hat\beta_1(x_i-\bar x)$), y
$\sum_i(x_i-\bar x)\hat u_i = \sum_i x_i\hat u_i - \bar x\sum_i\hat u_i
= 0-0=0$ por las dos ecuaciones normales. Luego
$\mathrm{SST}=\mathrm{SSE}+\mathrm{SSR}$.

(b) $\hat y_i-\bar y=\hat\beta_1(x_i-\bar x)$, así que
$\mathrm{SSE}=\hat\beta_1^2\sum_i(x_i-\bar x)^2 = \hat\beta_1^2 S_{xx}
= (S_{xy}/S_{xx})^2 S_{xx} = S_{xy}^2/S_{xx}$.

(c) $R^2 = \dfrac{\mathrm{SSE}}{\mathrm{SST}}
= \dfrac{S_{xy}^2/S_{xx}}{S_{yy}} = \dfrac{S_{xy}^2}{S_{xx}S_{yy}}
= \hat\rho_{xy}^2$.
\end{proof}

\begin{exercise}[Regresión de $y$ sobre $x$ vs.\ de $x$ sobre $y$]\label{ej:reversa}
Sean $\hat\beta_{y|x}=S_{xy}/S_{xx}$ la pendiente de regresar $y$ sobre
$x$, y $\hat\beta_{x|y}=S_{xy}/S_{yy}$ la de regresar $x$ sobre $y$ (los
papeles invertidos). Demuestre que
\[
  \hat\beta_{y|x}\cdot\hat\beta_{x|y} = \hat\rho_{xy}^2 \le 1,
\]
con igualdad si y sólo si $|\hat\rho_{xy}|=1$ (correlación muestral
perfecta). Explique por qué esto muestra que, en general, la recta de
regresar $y$ sobre $x$ \textbf{no} es la misma que la de regresar $x$
sobre $y$ ``despejada'' para $y$.
\end{exercise}
\begin{proof}[Solución]
$\hat\beta_{y|x}\hat\beta_{x|y} = \dfrac{S_{xy}}{S_{xx}}\cdot
\dfrac{S_{xy}}{S_{yy}} = \dfrac{S_{xy}^2}{S_{xx}S_{yy}} =
\hat\rho_{xy}^2$, y $\hat\rho_{xy}^2\le1$ por la versión muestral de la
desigualdad de Cauchy--Schwarz, con igualdad sólo si $|\hat\rho_{xy}|=1$.
Si despejar $x$ de $\hat y=\hat\beta_0+\hat\beta_{y|x}x$ diera la recta
de regresar $x$ sobre $y$, su pendiente (en $y$) sería
$1/\hat\beta_{y|x}$; pero la pendiente real de esa segunda regresión es
$\hat\beta_{x|y}$, y $1/\hat\beta_{y|x}=\hat\beta_{x|y}$ sólo cuando
$\hat\beta_{y|x}\hat\beta_{x|y}=1$, es decir, sólo en el caso extremo
$|\hat\rho_{xy}|=1$. Fuera de ese caso, las dos rectas de regresión son
genuinamente distintas.
\end{proof}

\begin{exercise}[Aplicación: beta de un activo]\label{ej:capm}
Un analista corre la regresión $R_i - r_f = \alpha + \beta(R_m-r_f)+u$
(la ecuación de mercado del CAPM) con datos mensuales y obtiene
$S_{xy}=0.0056$, $S_{xx}=0.0080$, $\overline{R_i-r_f}=0.012$,
$\overline{R_m-r_f}=0.010$. Calcule $\hat\beta$ y $\hat\alpha$, e
interprete ambos en el lenguaje del CAPM (¿el activo es más o menos
volátil que el mercado? ¿hay evidencia de rendimiento anormal?).
\end{exercise}
\begin{proof}[Solución]
$\hat\beta = 0.0056/0.0080 = 0.70$. Como $\hat\beta<1$, el activo es
--según esta muestra-- menos volátil que el mercado (un movimiento de
1\% en el exceso de mercado se asocia, en promedio, con un movimiento de
sólo 0.70\% en el exceso del activo). El intercepto es
$\hat\alpha = 0.012-0.70(0.010)=0.012-0.007=0.005$, es decir, un
``alfa'' positivo de 0.5\% mensual no explicado por el riesgo de
mercado --potencial evidencia de rendimiento anormal, aunque una
conclusión firme requeriría una prueba de hipótesis sobre $\hat\alpha$
(Sección 3 del syllabus), no sólo su signo.
\end{proof}

\begin{exercise}[MCO con un regresor binario: diferencia de medias]\label{ej:dummy}
Suponga que $x_i\in\{0,1\}$ es una variable \emph{dummy} (por ejemplo,
$x_i=1$ si la empresa $i$ paga dividendos, $0$ si no), con $n_1$
observaciones en el grupo $x=1$ y $n_0=n-n_1$ en el grupo $x=0$. Sean
$\bar y_1$ y $\bar y_0$ las medias muestrales de $y$ dentro de cada
grupo. Demuestre que
\[
  \hat\beta_1 = \bar y_1 - \bar y_0, \qquad \hat\beta_0 = \bar y_0 .
\]
\end{exercise}
\begin{proof}[Solución]
Como $x_i\in\{0,1\}$, $\bar x=n_1/n$, y
$\bar y = (n_0\bar y_0+n_1\bar y_1)/n$. Para las observaciones con
$x_i=1$, $x_i-\bar x = 1-n_1/n=n_0/n$; para las de $x_i=0$,
$x_i-\bar x=-n_1/n$. Separando la suma $S_{xy}$ por grupo,
\begin{align*}
  S_{xy} &= \frac{n_0}{n}\sum_{i:\,x_i=1}(y_i-\bar y)
          - \frac{n_1}{n}\sum_{i:\,x_i=0}(y_i-\bar y) \\
  &= \frac{n_0}{n}\,n_1(\bar y_1-\bar y) - \frac{n_1}{n}\,n_0(\bar y_0-\bar y)
   = \frac{n_0n_1}{n}\big[(\bar y_1-\bar y)-(\bar y_0-\bar y)\big]
   = \frac{n_0n_1}{n}(\bar y_1-\bar y_0).
\end{align*}
Análogamente,
$S_{xx} = n_1\big(\tfrac{n_0}{n}\big)^2+n_0\big(\tfrac{n_1}{n}\big)^2
= \dfrac{n_0n_1(n_0+n_1)}{n^2} = \dfrac{n_0n_1}{n}$. Por lo tanto
\[
  \hat\beta_1 = \frac{S_{xy}}{S_{xx}}
  = \frac{(n_0n_1/n)(\bar y_1-\bar y_0)}{n_0n_1/n} = \bar y_1-\bar y_0,
\]
y $\hat\beta_0=\bar y-\hat\beta_1\bar x$. Sustituyendo
$\bar y=(n_0\bar y_0+n_1\bar y_1)/n$ y $\bar x=n_1/n$:
\[
  \hat\beta_0 = \frac{n_0\bar y_0+n_1\bar y_1}{n}
  - (\bar y_1-\bar y_0)\frac{n_1}{n}
  = \frac{n_0\bar y_0+n_1\bar y_1 - n_1\bar y_1+n_1\bar y_0}{n}
  = \frac{(n_0+n_1)\bar y_0}{n} = \bar y_0 . \qedhere
\]
\end{proof}

\begin{summaryBox}
Regresar $y$ sobre una dummy es, algebraicamente, exactamente lo mismo
que comparar la media de $y$ entre dos grupos: $\hat\beta_0$ es la media
del grupo base ($x=0$) y $\hat\beta_1$ es la diferencia de medias entre
grupos. Ésta es la base de por qué, más adelante en el curso (Tema 4,
variables cualitativas), MCO con variables dummy se usa para probar
diferencias entre grupos con toda la maquinaria de inferencia de la
regresión (Tema 3): una prueba $t$ sobre $\hat\beta_1$ es, en este caso
particular, una prueba de diferencia de medias.
\end{summaryBox}

% ================================================================
\section{Estilo examen CFA (opción múltiple)}
% ================================================================

\begin{ideaBox}
Los exámenes CFA presentan la pendiente de MCO como
$\hat\beta_1=\Cov(x,y)/\Var(x)$, usando la covarianza y la varianza
\emph{muestrales estándar} (divididas entre $n-1$), en vez de las sumas
sin dividir $S_{xy},S_{xx}$ usadas en el resto de esta hoja. El primer
ejercicio muestra que ambas notaciones dan exactamente el mismo
número; los siguientes son preguntas de opción múltiple con el formato
de tres alternativas (A/B/C) típico del examen.
\end{ideaBox}

\begin{exercise}[De $S_{xy}/S_{xx}$ a $\Cov(x,y)/\Var(x)$]\label{ej:cfa-notacion}
La covarianza y varianza muestrales \emph{estándar} son
$\Cov(x,y)=S_{xy}/(n-1)$ y $\Var(x)=S_{xx}/(n-1)$. Demuestre que
$\Cov(x,y)/\Var(x) = S_{xy}/S_{xx} = \hat\beta_1$, y explique por qué no
importa si el software (o el examen) usa $n-1$ o $n$ en el denominador.
\end{exercise}
\begin{proof}[Solución]
\[
  \frac{\Cov(x,y)}{\Var(x)} = \frac{S_{xy}/(n-1)}{S_{xx}/(n-1)}
  = \frac{S_{xy}}{S_{xx}} = \hat\beta_1,
\]
porque el factor $1/(n-1)$ aparece igual en numerador y denominador y se
cancela. El mismo argumento funciona con cualquier otro divisor común
(por ejemplo $n$): $\hat\beta_1$ es un \emph{cociente} de dos cantidades
que se reescalan juntas, así que su valor no depende de la convención de
normalización que use un paquete estadístico en particular --sólo debe
usarse la \textbf{misma} convención en la covarianza y en la varianza a
la vez.
\end{proof}

\begin{exercise}[Pendiente a partir de covarianza y varianza reportadas]\label{ej:cfa-pendiente}
Un analista de renta variable estima
$R_{i,t}=\alpha+\beta R_{m,t}+u_t$ con 36 meses de datos y reporta
$\Cov(R_i,R_m)=0.00186$ y $\Var(R_m)=0.00150$ (estadísticas estándar,
ya divididas entre $n-1=35$). ¿Cuál de las siguientes opciones está más
cerca del valor de $\hat\beta$?
\begin{enumerate}[label=(\Alph*),nosep,leftmargin=2.2em]
  \item $0.81$
  \item $1.24$
  \item $1.55$
\end{enumerate}
\end{exercise}
\begin{proof}[Solución]
$\hat\beta=\Cov(R_i,R_m)/\Var(R_m) = 0.00186/0.00150 = 1.24$. Respuesta
(B). (Que $\Cov$ y $\Var$ ya estén divididas entre $n-1$ no cambia nada,
por el Ejercicio~\ref{ej:cfa-notacion}.)
\end{proof}

\begin{exercise}[Correlación y proporción de varianza no explicada]\label{ej:cfa-r2}
La correlación muestral entre el rendimiento de un fondo y el de su
benchmark es $\hat\rho=0.72$. Un cliente pregunta qué porcentaje de la
variación del rendimiento del fondo \textbf{no} queda explicado por el
benchmark. ¿Cuál de las siguientes opciones es la más cercana?
\begin{enumerate}[label=(\Alph*),nosep,leftmargin=2.2em]
  \item $28.0\%$
  \item $48.2\%$
  \item $51.8\%$
\end{enumerate}
\end{exercise}
\begin{proof}[Solución]
$R^2=\hat\rho^2=0.72^2=0.5184$ es la proporción \emph{explicada}
(51.8\%, opción (C) --una trampa clásica para quien confunde $R^2$ con
la proporción no explicada). Lo no explicado es
$1-R^2=1-0.5184=0.4816\approx48.2\%$: respuesta (B). (La opción (A),
$28\%$, es la trampa de calcular $1-\hat\rho$ en vez de $1-\hat\rho^2$
--olvidar elevar al cuadrado la correlación antes de restar de 1.)
\end{proof}

\begin{exercise}[Pronóstico puntual con la recta estimada]\label{ej:cfa-pronostico}
Con el modelo estimado en el Ejercicio~\ref{ej:numerico-grande},
$\hat y = 2.2+0.6143\,x$, ¿cuál es el valor pronosticado de $y$ para
$x=9$, el más cercano a...
\begin{enumerate}[label=(\Alph*),nosep,leftmargin=2.2em]
  \item $7.7$
  \item $5.5$
  \item $8.9$
\end{enumerate}
\end{exercise}
\begin{proof}[Solución]
$\hat y = 2.2+0.6143(9) = 2.2+5.5287 = 7.7287\approx7.7$. Respuesta (A).
(La opción (B), $5.5$, es la trampa de reportar sólo
$0.6143\times9\approx5.5$, olvidando sumar el intercepto.)
\end{proof}

\begin{exercise}[Sensibilidad económica: crecimiento de ventas y PIB]\label{ej:cfa-sensibilidad}
Un analista regresa el crecimiento anual de las ventas de una empresa
($y$, en \%) sobre el crecimiento del PIB ($x$, en \%) y obtiene
$\hat\beta_0=-0.3$, $\hat\beta_1=1.8$. Si el consenso de mercado
pronostica un crecimiento del PIB de $2.5\%$ el próximo año, ¿cuál de
las siguientes opciones está más cerca del crecimiento de ventas
pronosticado?
\begin{enumerate}[label=(\Alph*),nosep,leftmargin=2.2em]
  \item $4.2\%$
  \item $4.5\%$
  \item $4.8\%$
\end{enumerate}
\end{exercise}
\begin{proof}[Solución]
$\hat y = -0.3+1.8(2.5) = -0.3+4.5 = 4.2\%$. Respuesta (A). (La opción
(B) olvida el intercepto por completo; la (C) le cambia el signo,
sumando $0.3$ en vez de restarlo.)
\end{proof}

\begin{exercise}[Dummy: rendimiento promedio del grupo tratado]\label{ej:cfa-dummy}
Un analista regresa el rendimiento anual de 50 fondos ($y$) sobre una
variable dummy $D=1$ si el fondo es de gestión activa y $D=0$ si es
pasivo (30 fondos activos, 20 pasivos), y obtiene $\hat\beta_0=4.5\%$,
$\hat\beta_1=-0.8\%$. ¿Cuál de las siguientes opciones es la más
cercana al rendimiento promedio de los fondos de gestión \textbf{activa}?
\begin{enumerate}[label=(\Alph*),nosep,leftmargin=2.2em]
  \item $4.5\%$
  \item $3.7\%$
  \item $5.3\%$
\end{enumerate}
\end{exercise}
\begin{proof}[Solución]
Por el Ejercicio~\ref{ej:dummy}, $\hat\beta_0=\bar y_0$ es la media del
grupo \emph{base} ($D=0$, fondos pasivos) y $\hat\beta_1=\bar
y_1-\bar y_0$ es la diferencia (activo menos pasivo). El rendimiento
promedio del grupo activo ($D=1$) es entonces
$\bar y_1=\hat\beta_0+\hat\beta_1=4.5\%+(-0.8\%)=3.7\%$. Respuesta (B).
(La opción (A), $4.5\%$, es la trampa de reportar $\hat\beta_0$
--el rendimiento del grupo \emph{pasivo}, no del activo.)
\end{proof}

% ================================================================
\section{Reto (opcional)}
% ================================================================

\begin{exercise}[Por qué MCO y no otro estimador lineal]\label{ej:blue}
Considere estimadores \emph{lineales} de la pendiente de la forma
$\tilde\beta_1=\sum_i c_i y_i$, con pesos $c_i$ que dependen sólo de los
$x_i$ (no de $y$). Para que $\tilde\beta_1$ sea \emph{insesgado} bajo
$y_i=\beta_0+\beta_1x_i+u_i$ con $\E[u_i]=0$, se requiere
\[
  \sum_i c_i = 0 \qquad \text{y} \qquad \sum_i c_i x_i = 1 .
\]
Suponiendo además que los errores son homocedásticos e incorrelacionados
($\Var(u_i)=\sigma^2$, $\Cov(u_i,u_j)=0$ para $i\neq j$), se tiene
$\Var(\tilde\beta_1)=\sigma^2\sum_i c_i^2$. Demuestre que, entre todos
los pesos $c_i$ que cumplen las dos restricciones de insesgamiento, la
elección de MCO, $c_i=(x_i-\bar x)/S_{xx}$, es la que \textbf{minimiza}
$\sum_i c_i^2$ --es decir, MCO es el estimador lineal insesgado de
mínima varianza (el teorema de Gauss--Markov, en el caso de la
pendiente simple).
\end{exercise}
\begin{hintBox}
Escriba $c_i = \dfrac{x_i-\bar x}{S_{xx}} + d_i$ para pesos arbitrarios
$d_i$ que cumplen las restricciones, y muestre primero que las dos
restricciones de insesgamiento fuerzan $\sum_i d_i=0$ y
$\sum_i d_i(x_i-\bar x)=0$ (recuerde $\sum_i(x_i-\bar x)=0$). Luego
expanda $\sum_ic_i^2=\sum_i\Big(\dfrac{x_i-\bar x}{S_{xx}}+d_i\Big)^2$ y
muestre que el término cruzado se anula usando esas dos identidades,
dejando $\sum_ic_i^2 = \dfrac{1}{S_{xx}} + \sum_id_i^2 \ge \dfrac1{S_{xx}}$,
con igualdad si y sólo si $d_i=0$ para toda $i$ --es decir, si y sólo si
$c_i$ son los pesos de MCO. (Este es, de nuevo, el mismo tipo de
argumento de ``completar el cuadrado'' del
Ejercicio~\ref{ej:sin-calculo}, ahora en el espacio de los pesos $c_i$
en vez de en $b_1$.)
\end{hintBox}

\end{document}
